\textbf{Phần 1. Trắc nghiệm lựa chọn}
\setcounterpageref{ex}{0}

\begin{ex} %[4D3V16]
Số gồm $6$ trăm triệu, $4$ chục nghìn và $2$ đơn vị viết là
\choice
{$640\,000\,002$}
{$604\,000\,002$}
{\True $600\,040\,002$}
{$600\,400\,002$}
\loigiai{
Số đó là $600\,040\,002$.\\
\textbf{Đáp án của câu hỏi là C}
}
\end{ex}

\begin{ex} %[4D3V16]
Chữ số 8 trong số $485\,600\,123$ thuộc hàng nào?
\choice
{\True Hàng chục triệu}
{Hàng triệu}
{Hàng trăm triệu}
{Hàng trăm nghìn}
\loigiai{
Chữ số 8 ở hàng chục triệu.\\
\textbf{Đáp án của câu hỏi là A}
}
\end{ex}

\begin{ex} %[4D3V16]
Làm tròn số $7\,480\,000$ đến hàng trăm nghìn ta được
\choice
{$7\,400\,000$}
{$7\,000\,000$}
{$8\,000\,000$}
{\True $7\,500\,000$}
\loigiai{
Chữ số hàng chục nghìn là 8 nên làm tròn lên $7\,500\,000$.\\
\textbf{Đáp án của câu hỏi là D}
}
\end{ex}

\begin{ex} %[4D3V16]
Số liền trước của số $10\,000\,000$ là
\choice
{$10\,000\,001$}
{\True $9\,999\,999$}
{$9\,999\,990$}
{$900\,000$}
\loigiai{
$10\,000\,000 - 1 = 9\,999\,999$.\\
\textbf{Đáp án của câu hỏi là B}
}
\end{ex}

\begin{ex} %[4D3V16]
Trong các số sau, số nào lớn nhất?
\choice
{$98\,765\,432$}
{$99\,999\,999$}
{\True $102\,345\,678$}
{$89\,000\,000$}
\loigiai{
$102\,345\,678$ có 9 chữ số, lớn nhất.\\
\textbf{Đáp án của câu hỏi là C}
}
\end{ex}

\begin{ex} %[4D3V16]
Lớp nghìn của số $456\,789\,012$ gồm các chữ số
\choice
{\True $7; 8; 9$}
{$4; 5; 6$}
{$0; 1; 2$}
{$4; 7; 0$}
\loigiai{
Lớp nghìn gồm 7, 8, 9.\\
\textbf{Đáp án của câu hỏi là A}
}
\end{ex}

\begin{ex} %[4D3V16]
Tìm chữ số $x$ biết $234,x56 < 234\,156$:
\choice
{$x = 1$}
{$x = 2$}
{$x = 3$}
{\True $x = 0$}
\loigiai{
$x = 0$ thỏa mãn $234\,056 < 234\,156$.\\
\textbf{Đáp án của câu hỏi là D}
}
\end{ex}

\begin{ex} %[4D3V16]
Tìm số hạng tiếp theo của dãy số: $1; 3; 6; 10; 15; \dots$
\choice
{$18$}
{\True $21$}
{$20$}
{$25$}
\loigiai{
Quy luật cộng tăng dần $+2, +3, +4, +5, +6 \Rightarrow 15 + 6 = 21$.\\
\textbf{Đáp án của câu hỏi là B}
}
\end{ex}

\begin{ex} %[4D3V16]
Một thư viện số có $12\,450\,000$ tài liệu. Làm tròn số tài liệu này đến hàng trăm nghìn ta được
\choice
{$12\,400\,000$}
{$12\,000\,000$}
{\True $12\,500\,000$}
{$13\,000\,000$}
\loigiai{
Hàng chục nghìn là 5 nên làm tròn thành $12\,500\,000$.\\
\textbf{Đáp án của câu hỏi là C}
}
\end{ex}

\begin{ex} %[4D3V16]
Tổng số chẵn bé nhất có 7 chữ số và số lẻ lớn nhất có 6 chữ số là
\choice
{\True $1\,999\,999$}
{$1\,999\,998$}
{$2\,000\,000$}
{$1\,099\,999$}
\loigiai{
$1\,000\,000 + 999\,999 = 1\,999\,999$.\\
\textbf{Đáp án của câu hỏi là A}
}
\end{ex}

\begin{ex} %[4D3V16]
Số $M = 40\,000\,000 + 5\,000\,000 + 60\,000 + 700 + 8$ viết gọn là
\choice
{$45\,600\,708$}
{$45\,607\,008$}
{$405\,060\,708$}
{\True $45\,060\,708$}
\loigiai{
$45\,060\,708$.\\
\textbf{Đáp án của câu hỏi là D}
}
\end{ex}

\begin{ex} %[4D3V16]
Tìm số tự nhiên $x$ có 7 chữ số biết rằng $x$ là số chẵn chia hết cho 5, các chữ số thuộc lớp triệu bằng 8, các chữ số thuộc lớp nghìn bằng 6 và các chữ số thuộc lớp đơn vị bằng 0.
\choice
{$8\,666\,000$}
{\True $8\,666\,000$}
{$88\,666\,000$}
{$8\,000\,660$}
\loigiai{
Số có 7 chữ số nên lớp triệu có 1 chữ số là 8. Số đó là $8\,666\,000$.\\
\textbf{Đáp án của câu hỏi là B}
}
\end{ex}

\textbf{Phần 2. Trắc nghiệm đúng sai}
\setcounterpageref{ex}{0}

\begin{ex} %[4D3V16]
Cho số $123\,456\,789$. Xét tính đúng sai:
\choiceTF
{\True Số này gồm 3 lớp: lớp triệu, lớp nghìn và lớp đơn vị.}
{\True Chữ số 1 thuộc hàng trăm triệu.}
{Chữ số 4 thuộc hàng trăm nghìn có giá trị $40\,000$.}
{\True Chữ số 4 thuộc hàng trăm nghìn có giá trị $400\,000$.}
\loigiai{
\begin{itemize}
    \item Gồm 3 lớp: Đúng.
    \item Chữ số 1 hàng trăm triệu: Đúng.
    \item Chữ số 4 có giá trị $400\,000 \neq 40\,000$. Ý 3 Sai, Ý 4 Đúng.
\end{itemize}
\begin{itemize}
    \item a) \textbf{Đúng}.
    \item b) \textbf{Đúng}.
    \item c) \textbf{Sai}.
    \item d) \textbf{Đúng}.
\end{itemize}
}
\end{ex}

\begin{ex} %[4D3V16]
Xét tính đúng sai về so sánh và làm tròn:
\choiceTF
{\True $45\,678\,000 > 45\,670\,800$.}
{\True Làm tròn $8\,960\,000$ đến hàng trăm nghìn được $9\,000\,000$.}
{Làm tròn $1\,240\,000$ đến hàng trăm nghìn được $1\,300\,000$.}
{\True Làm tròn $1\,240\,000$ đến hàng trăm nghìn được $1\,200\,000$.}
\loigiai{
\begin{itemize}
    \item $45\,678\,000 > 45\,670\,800$: Đúng.
    \item $8\,960\,000 \rightarrow 9\,000\,000$: Đúng.
    \item $1\,240\,000 \rightarrow 1\,200\,000 \neq 1\,300\,000$. Ý 3 Sai, Ý 4 Đúng.
\end{itemize}
\begin{itemize}
    \item a) \textbf{Đúng}.
    \item b) \textbf{Đúng}.
    \item c) \textbf{Sai}.
    \item d) \textbf{Đúng}.
\end{itemize}
}
\end{ex}

\textbf{Phần 3. Trả lời ngắn}
\setcounterpageref{ex}{0}

\begin{ex} %[4D3V16]
Viết số: Chín mươi lăm triệu sáu trăm nghìn.
\shortans{95600000}
\loigiai{
\textbf{Đáp số:} $95600000$\\
$95\,600\,000$.
}
\end{ex}

\begin{ex} %[4D3V16]
Làm tròn số $4\,750\,000$ đến hàng trăm nghìn.
\shortans{4800000}
\loigiai{
\textbf{Đáp số:} $4800000$\\
$4\,800\,000$.
}
\end{ex}

\begin{ex} %[4D3V16]
Tìm số chẵn lớn nhất có 7 chữ số.
\shortans{9999998}
\loigiai{
\textbf{Đáp số:} $9999998$\\
$9\,999\,998$.
}
\end{ex}

\begin{ex} %[4D3V16]
Trong số $345\,678\,901$, tính tổng các chữ số thuộc lớp nghìn.
\shortans{21}
\loigiai{
\textbf{Đáp số:} $21$\\
$6 + 7 + 8 = 21$.
}
\end{ex}

\begin{ex} %[4D3V16]
Dãy số $1; 4; 7; 10; \dots; 31$ có bao nhiêu số hạng?
\shortans{11}
\loigiai{
\textbf{Đáp số:} $11$\\
$(31 - 1) : 3 + 1 = 11$.
}
\end{ex}

\begin{ex} %[4D3V16]
Tìm hiệu giữa số lớn nhất có 8 chữ số khác nhau và số bé nhất có 8 chữ số khác nhau.
\shortans{88530864}
\loigiai{
\textbf{Đáp số:} $88530864$\\
$98\,765\,432 - 10\,234\,568 = 88\,530\,864$.
}
\end{ex}
